\documentclass{beamer}
\usepackage[orientation=portrait,size=a0,scale=1.25]{beamerposter}

\usepackage[utf8]{inputenc}
\usepackage[T1]{fontenc}
\usepackage[english]{babel}
\usepackage{amsmath,amssymb}
\usepackage{graphicx}
\usepackage{booktabs}
\usepackage{multicol}
\usepackage{ragged2e}
\usepackage{lipsum}
\usepackage{tikz}
\usepackage{apacite}

\usetheme{simpleposter}

\graphicspath{{img/}}
\usetikzlibrary{arrows.meta,positioning}

\definecolor{posteraccent}{RGB}{19,61,148}

\renewcommand{\posterfooter}{\textbf{example.org/your-lab} \quad
first.last@example.org}

\title{A Concise Poster Title That Fits on Two Lines at Most}
\author{\LARGE First A. Author$^1$, Second B. Author$^{1,2}$, Third
C. Author$^3$}
\institute{
  $^1$Department, Institution One, City, Country. \\
  $^2$Laboratory, Institution Two, City, Country. \\
  $^3$Institution Three, City, Country.
}
\date{\today}

\begin{document}
\begin{frame}[t]{Keywords: \bfseries keyword one, keyword two, keyword three}
  \justifying

  \begin{columns}[t]

    \column{.48\textwidth}

    \begin{block}{Introduction}
      \lipsum[1]
    \end{block}

    \begin{block}{Motivation}
      \begin{columns}[t]
        \column{0.52\textwidth}
        \lipsum[2][1-5]
        \begin{itemize}
          \item First placeholder bullet point.
          \item Second placeholder bullet point.
          \item Third placeholder bullet point.
        \end{itemize}
        \column{0.44\textwidth}
        \begin{figure}
          \centering
          \placeholderbox{\linewidth}{15em}{FIGURE}
          \caption{Replace with \texttt{\textbackslash
          includegraphics} from \texttt{img/}.}
          \label{fig:placeholder}
        \end{figure}
      \end{columns}
    \end{block}

    \begin{block}{Methods}
      \lipsum[3][1-6]

      \begin{equation}
        \hat{\theta} = \arg\min_{\theta} \; \sum_{i=1}^{n}
        \ell\big(y_i, f_\theta(x_i)\big) + \lambda \lVert \theta \rVert_2^2
        \label{eq:objective}
      \end{equation}

      Equation~\eqref{eq:objective} is a placeholder; $\lambda$
      controls regularisation.

      \begin{figure}
        \centering
        \resizebox{\linewidth}{!}{%
          \begin{tikzpicture}[
              node distance=6em,
              every node/.style={font=\Large},
              box/.style={draw=posteraccent,fill=posteraccent!10,line width=2pt,
                rounded corners,minimum width=11em,minimum
              height=5em,align=center},
              arr/.style={-{Stealth[length=5mm]},line
            width=2pt,draw=posteraccent}]
            \node[box] (a) {Input};
            \node[box,right=of a] (b) {Process};
            \node[box,right=of b] (c) {Model};
            \node[box,below=4em of c] (d) {Output};
            \draw[arr] (a) -- (b);
            \draw[arr] (b) -- (c);
            \draw[arr] (c) -- (d);
            \draw[arr] (d.west) -| node[pos=0.25,above]{feedback} (b.south);
        \end{tikzpicture}}
        \caption{Placeholder pipeline drawn in Ti\emph{k}Z.}
        \label{fig:pipeline}
      \end{figure}
    \end{block}

    \column{.48\textwidth}

    \begin{block}{Results}
      \lipsum[4][1-5]

      \begin{table}
        \centering
        \begin{tabular}{lrrr}
          \toprule
          Method & Metric A & Metric B & Metric C \\
          \midrule
          Baseline    & 0.00 & 0.00 & 0.00 \\
          Variant I   & 0.00 & 0.00 & 0.00 \\
          Variant II  & 0.00 & 0.00 & 0.00 \\
          \textbf{Ours} & \textbf{0.00} & \textbf{0.00} & \textbf{0.00} \\
          \bottomrule
        \end{tabular}
        \caption{Placeholder results table.}
        \label{tab:results}
      \end{table}

      \begin{columns}[t]
        \column{0.48\textwidth}
        \placeholderbox{\linewidth}{14em}{PLOT A}
        \column{0.48\textwidth}
        \placeholderbox{\linewidth}{14em}{PLOT B}
      \end{columns}
    \end{block}

    \begin{block}{Discussion}
      \lipsum[5][1-5]
      Reference placeholder: \cite{placeholder2020,placeholder2021}.
    \end{block}

    \begin{block}{Conclusions \& Future Work}
      \begin{itemize}
        \item Placeholder takeaway one.
        \item Placeholder takeaway two.
        \item Placeholder next step.
      \end{itemize}
    \end{block}

    \begin{block}{Acknowledgements}
      \lipsum[6][1-3]
    \end{block}

    \begin{block}{References}
      \small
      \begin{multicols}{2}
        \setbeamertemplate{bibliography item}[text]
        \bibliographystyle{apacite}
        \bibliography{refs}
      \end{multicols}
    \end{block}

  \end{columns}
\end{frame}
\end{document}
